\section{Praktikum}
\label{sec:praktikum-bab6}

\begin{praktikumbox}
\vspace{1.1em}
\textbf{Judul:} Menghitung dan memverifikasi alamat fisik dengan
\eng{debugger}.\\[0.3em]
\textbf{Sub-CPMK:} 1.4 \quad\textbullet\quad
\textbf{Alokasi waktu:} 100 menit \quad\textbullet\quad
\textbf{Bentuk:} individu
\end{praktikumbox}

\subsection{Tujuan Praktikum}

Setelah menyelesaikan praktikum ini, mahasiswa mampu menghitung
alamat fisik dari pasangan \eng{segment:offset}, membuktikan hasilnya
dengan DEBUG dan emu8086, serta menunjukkan \eng{overlapping} segmen.

\subsection{Alat dan Bahan}

\begin{itemize}
    \item Komputer dengan emu8086 terpasang.
    \item DOSBox dengan TASM, TLINK, dan DEBUG.
    \item Berkas \file{alamat-fisik.asm} dan \file{hitung-alamat.asm}
          dari \file{code/asm/}.
    \item Kalkulator heksadesimal atau lembar kerja perhitungan.
\end{itemize}

\subsection{Langkah Kerja}

\begin{enumerate}
    \item \textbf{Merakit dan menjalankan program pertama.}
    \begin{enumerate}
        \item Buka DOSBox, lalu masuk ke direktori kerja.
        \item Rakit \file{alamat-fisik.asm} dengan
              \file{tasm alamat-fisik.asm}, lalu link dengan
              \file{tlink alamat-fisik.obj}.
        \item Jalankan program dan catat kedua nilai offset yang
              tampil.
    \end{enumerate}

    \item \textbf{Menghitung alamat fisik secara manual.}
    \begin{enumerate}
        \item Jalankan \file{debug alamat-fisik.exe}.
        \item Ketik \file{r}, lalu catat nilai \reg{DS}.
        \item Hitung alamat fisik setiap variabel dengan rumus
              $(\text{segmen} \times \hex{10}) + \text{offset}$.
              Tuliskan langkahnya, bukan hanya hasil akhir.
        \item Susun hasil pada tabel seperti
              Tabel~\ref{tab:analisis-offset}.
    \end{enumerate}

    \item \textbf{Memverifikasi dengan DEBUG.}
    \begin{enumerate}
        \item Tampilkan isi memori dengan
              \lstasm|D DS:0000 L 08|.
        \item Bandingkan isi yang tampil dengan nilai variabel pada
              program: \code{var\_byte} bernilai \hex{10} dan
              \code{var\_word} bernilai \hex{1234}.
        \item Perhatikan urutan byte pada \code{var\_word}. Jelaskan
              mengapa byte rendah tampil lebih dahulu.
    \end{enumerate}

    \item \textbf{Membuktikan overlapping.}
    \begin{enumerate}
        \item Bila \reg{DS} $= \hex{076A}$ dan offset
              \code{var\_word} $= \hex{0001}$, alamat fisiknya adalah
              \hex{076A1}.
        \item Tuliskan pasangan lain yang menunjuk alamat tersebut,
              misalnya \code{0769h:0011h}.
        \item Tampilkan kedua pasangan dengan perintah
              \lstasm|D| dan pastikan isinya sama.
        \item Jelaskan mengapa dua pasangan berbeda dapat menunjuk
              lokasi yang sama.
    \end{enumerate}

    \item \textbf{Menjalankan program kedua.}
    \begin{enumerate}
        \item Rakit dan jalankan \file{hitung-alamat.asm}.
        \item Bandingkan alamat fisik yang dihitung program dengan
              hasil perhitungan manual.
        \item Tambahkan satu variabel baru di depan \code{var\_nilai},
              rakit ulang, lalu jelaskan perubahan offset yang
              terjadi.
    \end{enumerate}

    \item \textbf{Kuis 1.}
    \begin{enumerate}
        \item Kerjakan kuis secara individu pada akhir pertemuan.
        \item Materi kuis mencakup perhitungan alamat fisik,
              paragraf 16 byte, \eng{overlapping} segmen, dan
              pembagian memori 1 MB.
    \end{enumerate}
\end{enumerate}

\subsection{Hasil yang Dikumpulkan}

\begin{enumerate}
    \item Lembar perhitungan manual alamat fisik untuk minimal enam
          pasangan \eng{segment:offset}.
    \item Tangkapan layar keluaran \file{alamat-fisik.asm} dan
          \file{hitung-alamat.asm}.
    \item Tangkapan layar sesi DEBUG dengan perintah \file{r} dan
          \file{d}.
    \item Bukti pengamatan \eng{overlapping} beserta penjelasan.
    \item Jawaban Kuis 1.
\end{enumerate}

\clearpage
\begin{asesmenbox}
\vspace{1.1em}
\textbf{Asesmen praktikum 4 (bobot 5\% dari komponen praktikum)} dan
\textbf{Kuis 1 (bobot 10\% dari komponen kuis)}.

\begin{tabular}{@{}p{6.4cm}cp{4.4cm}@{}}
    \toprule
    \textbf{Indikator} & \textbf{Bobot} & \textbf{Kriteria} \\
    \midrule
    Ketepatan perhitungan alamat fisik & 35\% &
    Seluruh langkah ditulis dan hasilnya benar \\
    \addlinespace
    Ketepatan verifikasi dengan DEBUG & 30\% &
    Perintah yang dipakai tepat dan isinya ditafsirkan dengan benar \\
    \addlinespace
    Analisis \eng{overlapping} & 20\% &
    Penjelasan pasangan alamat dan alamat fisik tepat \\
    \addlinespace
    Kerapian dokumentasi & 15\% &
    Langkah kerja terdokumentasi dan dapat diulang \\
    \bottomrule
\end{tabular}
\end{asesmenbox}
